% !TEX root = main.tex
\pagebreak
\section{Discussion}
\label{sec:discussion}

\subsection{Architectural Stability}

Our analysis of the Ceph file system at both the \ac{RADOS} and \ac{CephFS} layers shows that data throughput scales well, particularly when the file system is subjected to high-concurrency workloads with large file accesses. Ceph's decentralized approach using \ac{CRUSH} ensures that an increasing number of clients distribute their workloads across all available \acp{OSD} on the system. However, the per-task performance using \ac{TCP} is low compared to the overall achievable throughput, which could pose a potential challenge for workloads which use a single-rank file access pattern. This per-node performance ceiling is consistent with findings from larger-scale deployments: benchmarks at Clyso~\cite{ceph-1tib} reported a similar per-node \ac{IOPS} limitation of 400--600K random read \ac{IOPS}, which they likely attributed to Ceph's \ac{OSD} threading model and messenger implementation: this persisted regardless of the underlying storage performance. However, provided that the file access patterns are sufficiently parallel, Ceph can be a capable system for \ac{HPC} scratch or archival storage with large block \ac{I/O} patterns.

While the data-path results were encouraging, metadata performance within Ceph lies in stark contrast to \ac{RADOS} with a complex performance profile. Metadata operations, such as file creation and deletion, directory listings, and index node queries, do not demonstrate the same strong scaling performance across runs in the same way that \ac{RADOS} does. Ceph's \ac{MDS} daemons are responsible for all metadata operations at the \ac{CephFS} layer, resulting in a performance bottleneck. While the number of active \ac{MDS} daemons may be increased by changing the max\_mds parameter within the Ceph configuration, the benchmarking data suggest that scaling is not as straightforward as adding \ac{MDS} daemons to the cluster---performance rarely scales linearly with the number of \ac{MDS} daemons, and care must be taken in order to optimize file and folder layouts such that they can be distributed across multiple \ac{MDS} daemons evenly. A large number of methods are available for partitioning of folder structures across \ac{MDS} daemons, including automatic subtree partitioning, but it may be necessary for system administrators to consider whether manual adjustment of the placement policy is necessary. In our testing with random subtree pinning, Ceph did not consistently demonstrate the expected performance uplift with an increased number of \ac{MDS} daemons, as some \ac{MDS} ranks received a significantly higher share of files and folders relative to others, resulting in inefficient utilization.

Further empirical evaluation is required in order to validate whether our observed throughput inconsistencies stem from the inherent limitations of the \ac{MDS}, especially in real-world conditions. Our specific approach of subtree pinning, necessitated by the nature of our synthetic benchmarks, could have resulted in the suboptimal allocation of directories to \ac{MDS} daemons. In addition, we did not test features such as automatic subtree rebalancing based on collected load information due to the short-lived nature of the benchmark runs leaving no time for the subtree balancer to run.

\subsubsection{Write Performance}

A notable finding from our large-scale benchmarks is the write performance limitation observed with replicated pools on the \ac{RADOS} layer. Write throughput reached 10--15\,GiB/s regardless of the number of client nodes, while read throughput continued to scale well with additional nodes. The bottleneck appears to be on the cluster side during writes. In a replica-3 configuration, each client write triggers two additional writes between \acp{OSD}~\cite{ceph-paper}, tripling the internal network traffic relative to reads. This write amplification likely results in saturation of network resources on the cluster side before the client nodes are fully utilized. However, the read-to-write performance ratio in our benchmarks generally exceeded 3:1 at higher node counts, where a 3:1 ratio would be expected from write amplification alone. This suggests additional overhead from write operations within the \acp{OSD}. The Clyso evaluation~\cite{ceph-1tib} observed a related issue at scale, where placement groups entered a ``laggy'' state under high write concurrency, further degrading write throughput---an issue attributed to messenger overhead or contention within the \ac{OSD}.

\subsubsection{Applying In-Memory Results to Production}

Our evaluation employed two complementary benchmarking environments: an in-memory deployment for high-performance testing on the \ac{HPC} cluster, which isolates Ceph's software overhead by eliminating storage media as a bottleneck, and a controlled lab deployment with disk \ac{I/O} throttling at 50\,MiB/s per \ac{OSD}, which allows us to observe Ceph's behaviour under well-defined constraints. By connecting the results from both environments, we can provide a more complete picture of expected real-world performance.

The in-memory benchmarks establish a performance ceiling: for example, sequential write throughput with replica-3 \ac{RADOS} pools plateaued at 10--15\,GiB/s regardless of the number of client nodes, and per-node read throughput saturated slightly above 100\ Gbit/s. These figures represent the maximum throughput achievable when storage media is not a limiting factor, and is instead bounded by Ceph's and the host's networking stack as well as internal \ac{OSD} processing overhead. In a production deployment, actual throughput will be the minimum of this software ceiling and the aggregate throughput of the storage media in use, which is especially relevant for slow storage media such as \acp{HDD}.

Our lab benchmarks provide a complementary data point: with disk performance throttled to a known level, Ceph achieved approximately 70--80\% of the theoretical aggregate disk throughput for replicated pools. Combining this disk efficiency factor with the software ceiling from our in-memory results, the expected throughput of a production deployment can be estimated. If the aggregate disk throughput, taking into consideration the overhead of the pool type (replication or \ac{EC}), falls below the software ceiling, the deployment will likely be disk-limited, and the 70--80\% efficiency we observed in our lab benchmarks provides an estimate of the achievable performance. However, this 70--80\% efficiency factor is derived from slow (50\,MiB/s) virtual disks, and cannot be safely generalized to fast physical media such as modern \ac{NVMe} \acp{SSD}, due to the possibility of \ac{CPU} bottlenecks within the BlueStore backend. If the aggregate disk throughput exceeds the software ceiling, as would be the case with high-performance network adapters and a large number of high-performance \acp{SSD}, the deployment will likely be software-limited: our in-memory results provide the upper bound which can be applied to this scenario. Furthermore, additional overhead introduced by BlueStore is not accounted for in our MemStore benchmarks, and may further reduce throughput in disk-limited scenarios.

\subsection{Network Performance}

Our experimental results highlight the importance of the role of high-performance networking in Ceph's storage performance. Ceph's reliance on the regular \ac{TCP/IP} stack introduces latency and throughput penalties, especially at small file sizes, which reduced overall performance for certain access patterns. Unlike incumbent \ac{HPC} file systems, Ceph does not provide a production-grade \ac{RDMA} transport: while the experimental \ac{RDMA} transport was tested in this thesis, we observed poor scaling results as well as the failure to achieve higher peak throughput relative to \acs{TCP} for large file accesses. However, small-file accesses as well as accesses with low numbers of clients proved significantly faster with \ac{RDMA}, likely due to its lower \ac{CPU} overhead and end-to-end latency characteristics.

Even though the \ac{RDMA} transport is available and may perform better in certain narrow edge cases, the experimental and undocumented nature of this feature is unlikely to be of use in current \ac{HPC} clusters due to its instability and uncertain future prospects with the introduction of the Crimson backend.

\subsection{\acs{CephFS} Overhead}

A direct comparison of the \ac{RADOS} and \ac{CephFS} results we gathered on the same cluster allows for us to quantify the overhead introduced by the \ac{CephFS} layer. For large sequential \ac{I/O} (64\,MiB), the \ac{CephFS} throughput scaled similarly to \ac{RADOS}, indicating that the \ac{CephFS} overhead is minimal for bulk data transfers where the \ac{MDS} is only involved at the beginning and end of the transfer for the file open/close operations. However, for small \ac{I/O} operations (4\,KiB), the overhead was substantial: at 8 nodes with 1 task per node, \ac{RADOS} achieved 45.9K read \ac{IOPS} compared to 1.65K for \ac{CephFS}---a 27$\times$ difference. This gap narrowed at higher concurrency levels, but \ac{CephFS} consistently remained an order of magnitude slower than \ac{RADOS} for small reads. Write performance showed a smaller overhead, with the \ac{MDS} becoming the primary bottleneck: a single \ac{MDS} produced approximately 6\ k\acs{IOPS}, while 4 \ac{MDS} daemons with distributed pinning achieved up to 11\ k\acs{IOPS}. These findings suggest that the \ac{CephFS} overhead is highly workload-dependent. For large sequential transfers the overhead is negligible, but potentially prohibitive for small-file, metadata-intensive workloads. For \ac{HPC} deployments, this implies that Ceph's suitability depends heavily on the dominant access pattern of the target workload---modern \ac{AI} workloads may be particularly metadata-intensive.

\subsection{Storage Efficiency}

In contrast to other storage systems, Ceph provides a flexible software-defined mechanism for ensuring the durability of data. Cluster administrators are not required to choose a fixed \ac{RAID} configuration. Lustre, when using ZFS as its underlying file system, is limited by the traditional considerations of \ac{RAID} configurations\footnote{\url{https://wiki.lustre.org/ZFS}. Accessed 26.02.2026.}. Since erasure-coding or replication can be defined at a pool level, different pools can be equipped with different policies: archival storage focusing on space efficiency over performance may choose to use a more efficient erasure-coding profile, while scratch storage may forgo the computationally expensive erasure coding in favour of simple replication---also on top of the same pool of physical disks and \acp{OSD} if desired.

\subsubsection{Erasure Coding}

The erasure-coding method presented by Ceph has several potential advantages for \ac{HPC} system administrators.

As erasure coding splits data into chunks, this results in, e.g., a 100\,MiB file being split into 4 data chunks and 2 coding chunks across 6 \acp{OSD} for a total of 150\,MiB written, or 25\,MiB per \ac{OSD}. A replica-3 profile, which simply creates 3 copies of data, results in at least 300 MiB of data written to the disk, by definition. In addition, 100\,MiB is written to per \ac{OSD} in such a configuration. 

If the \ac{OSD} is bandwidth-limited to 50\,MiB/s, the data would take at least 2 seconds to be written to the replica-3 \acp{OSD}, while the erasure-coded configuration could potentially complete the write in 0.5 seconds.

However, if the file sizes are relatively small, the latency penalty of erasure coding may be significant. However, with our memory-backed \acp{OSD}, we did not observe a drastic change in performance characteristics. Additionally, if large amounts of data are being written, the \ac{CPU} overhead could potentially outweigh the performance benefits: this was apparent in the case study of the 1\ TiB/s Ceph cluster~\cite{ceph-1tib}, although the performance optimizations mentioned within it seem to have been merged since this report\footnote{\url{https://github.com/ceph/ceph/pull/55196}. Accessed 26.02.2026.}.

Our controlled disk performance tests also corroborates the theoretical advantages of erasure coding over replication. \ac{EC} pools consistently outperformed replicated pools in our lab conditions, where disk throughput was the major limiting factor. The striping behaviour of \ac{EC} improved the performance of single-object reads and writes substantially, especially when the number of chunks were high, and resulted in much better write performance due to the \ac{CPU} overhead being negligible at the 50\,MiB/s disk performance we used in our tests. Additionally, the reduced overall write load reduces the total network traffic as well as the total amount of data which must be written to disks, potentially making network utilization more efficient, as well as preserving the limited lifespan of physical storage media. However, this assumes that the overhead of erasure coding and potential journaling is insignificant. Our simulated production cluster's conditions show that for slow disks, the overhead of \ac{EC} can theoretically be outweighed by the performance benefits, but we were unable to quantify this for our purely in-memory cluster used in the large-scale benchmarks.

\subsection{Operational Considerations}

From a system administrator's perspective, Ceph exhibits a high degree of flexibility compared to traditional \ac{HPC} file systems. Our limited testing experience has shown that Ceph is user-friendly to set up, and uses technologies such as containerization in order to decouple itself from the software stack of the host system, further increasing operating system compatibility. All Ceph daemons within a production deployment with cephadm run within a container runtime, such as Docker or Podman, providing a clear administrative advantage over Lustre: because Lustre relies on highly customized, out-of-tree Linux kernel modules for its system, \ac{HPC} administrators may be locked into specific, older operating systems. In Lustre's case specifically, the current public version at the time of this thesis\footnote{Version 2.17.0, as of 25.02.2026.} provides kernel modules for \ac{RHEL} 10, \ac{SLES} 15, and Ubuntu 24.04 for the client, and server modules only for \ac{RHEL} 9\footnote{\url{https://downloads.whamcloud.com/public/lustre/}. Accessed 26.02.2026.}. The tight coupling of the storage system to the kernel does, at least in part, allow Lustre to achieve its high performance results on the IO500~\cite{io500-full-list}, but it complicates system upgrades and the use of heterogeneous software installations across multiple systems. Ceph bypasses these limitations entirely by running the server component entirely within containers and using standard \ac{TCP/IP} networking, while clients benefit from in-tree kernel module support and \ac{FUSE} for kernels without---or with an older version of---\ac{CephFS}. However, Lustre aims to transition to an upstream kernel module as well\footnote{\url{https://wiki.lustre.org/Lustre_Upstreaming_to_Linux_Kernel}. Accessed 26.02.2026.}, which could make it more competitive with Ceph in this regard.

Additionally, the widespread adoption of Ceph further promotes its long-term viability as a file system suitable for \ac{HPC}. In addition to specific \ac{HPC} users such as CERN~\cite{cern-evaluating-cephfs-perf}, major cloud and hosting providers such as DigitalOcean\footnote{\url{https://www.digitalocean.com/blog/why-we-chose-ceph-to-build-block-storage}. Accessed 26.02.2026.} and Hetzner\footnote{\url{https://docs.hetzner.com/storage/general/which-storage-is-right-for-me}. Accessed 26.02.2026.} use the technology, with vendors such as Red Hat\footnote{\url{https://www.redhat.com/en/technologies/storage/ceph}. Accessed 26.02.2026.} and Clyso\footnote{\url{https://www.clyso.com/eu/en/products/clyso-enterprise-storage/}. Accessed 26.02.2026.} providing commercial offerings for Ceph.

Beyond its file system layer, Ceph also supports multiple other services, such as \ac{RBD} for block storage, used by the aforementioned cloud providers for \ac{VM} block storage, and \ac{RGW} for \ac{S3} compatible object storage, enabling mixed uses for the same underlying infrastructure. A single Ceph cluster could be used to serve \ac{HPC} scratch storage, home directories for users, block storage for virtual machines, and object storage for data analytics. Such convergence could improve hardware utilization, which could potentially reduce storage costs.

In terms of system administration, our evaluation of the cephadm tool and the Ceph dashboard was generally positive. The Ceph dashboard provides a simple overview of the cluster, and allows administrators to perform basic management tasks without specific expertise in the Ceph \ac{CLI}, while its underlying Prometheus-based monitoring stack further enables integration with external tools such as Grafana. 

However, we also observed issues with the consistency of data reported by the administration tools, leading to some confusion when performing management tasks via the \ac{CLI}. Additionally, the reinstallation of the cluster required manual intervention in order to restore the systems to a clean state for a fresh installation. While none of our observed issues are critical, Ceph's administrative tooling could be a potential source of improvement for future versions.

\subsection{Limitations}

\subsubsection{Infrastructure}
\label{sec:infrastructure-limitations}

Our empirical evaluation of Ceph was limited by several infrastructure limitations. Our benchmarks were primarily executed in a rootless environment with fully in-memory storage backends (MemStore \acp{OSD}) rather than writing to physical storage devices. While we consider the MemStore approach to be effective at eliminating compounding factors, isolating storage performance from the observed performance overhead, we were unable to evaluate the characteristics of hard drives and flash storage as it relates to Ceph; real-world deployments contend with access latency and significantly lower throughput limitations, and the impact of flushing data to disk is not clear from our report.

Consequently, our results gathered with MemStore should be interpreted as an upper bound on Ceph's achievable performance for a given cluster configuration, where storage media latency and throughput limitations are not limiting factors. Production deployments with physical storage will be bounded by this software-derived ceiling in addition to the characteristics of the storage devices themselves. Interaction between Ceph's BlueStore backend and physical storage media, such as the write-ahead log and disk scheduling overhead, introduce additional performance overhead absent from our MemStore benchmarks. A full performance model for real-world Ceph performance would require the validation of the model against physical storage performance,  which we identify as an important direction for future work.

\subsubsection{Networking}

Furthermore, the limitations of the Omni-Path network combined with our specific version of Rocky Linux 8 containing a kernel bug resulted in intermittent and unpredictable node failures, preventing full-scale tests with the Omni-Path fabric. Therefore, our benchmarks were limited to a total of 16 combined client and server nodes, rather than the 256 or more which would have been theoretically possible on the Omni-Path equipped compute nodes. Therefore, we were unable to test the performance characteristics at a production scale, and the limitation in node counts made it difficult to definitively isolate whether our observed performance limitations were a result of network or \ac{CPU} bottlenecks on the client and server side---especially for our large-object throughput tests.

\subsubsection{IO500 Benchmark}

The lack of dedicated physical disk and low number of nodes prevented the execution of the standard IO500 benchmark suite under its official testing parameters. By definition, as our cluster is an in-memory cluster, it is not possible to run the IO500 benchmark suite under its official testing parameters\footnote{\url{https://io500.org/rules/submission}, rule 6: \enquote{All data must be written to persistent storage within the measured time for the individual benchmark, e.g. if a file system caches data, it must ensure that data is persistently stored before acknowledging the close.}. Accessed 01.03.2026.}. Additionally, the number of changes which would have to be made to the benchmarking options in order to allow it to run on our hardware would have been significant, and would have hindered the validity and comparability of the results. Our test cluster was equipped with 8\ TiB (16 $\times$ 512\,GiB) of \ac{RAM}, of which ~6\ TiB could be allocated safely to \ac{OSD} storage, which is insufficient for running the IO500 benchmark.

\subsubsection{Methodology}

Our evaluation scope did not include the simulation of complex network or hardware failures in production systems. Prior work conducted by CERN~\cite{cern-evaluating-cephfs-perf} demonstrated issues such as high-latency \acp{OSD} impacting throughput significantly due to hardware issues. However, we could not observe such performance degradation on our cluster---the uniform hardware profile and in-memory \acp{OSD} prevents observation of such phenomena. Additionally, because the cluster was an ephemeral test cluster and did not host production workloads, we were unable to observe the impact of Ceph's optimization features, such as the metadata and \ac{CRUSH} map rebalancing systems, as these are impacted by user workloads and run over long periods of time.

\pagebreak
\section{Future Work}
\label{sec:future-work}

\subsection{Upcoming Storage Backends}

The most interesting Ceph development for future research is Ceph's upcoming Crimson/SeaStore backend update\footnote{\url{https://ceph.io/en/news/crimson/}. Accessed 26.02.2026.}. As our work has demonstrated, the performance of Ceph's networking stack can be a limiting factor to performance, resulting in increased latency and reduced throughput compared to theoretical limits. The Crimson backend promises a complete overhaul of the storage and networking backend within Ceph, utilizing technologies such as \ac{DPDK} and \ac{SPDK}, supporting features such as kernel-bypass and reduced \ac{CPU} usage---features from which incumbent \ac{HPC} storage platforms currently benefit, although in different ways. Currently, the SeaStore backend is not suitable for production use, and according to the Ceph documentation is currently missing critical features, such as erasure coding and \ac{OSD} MClock for fair \ac{I/O} scheduling.

\subsection{Physical Storage Validation}

As discussed in~\cref{sec:infrastructure-limitations}, our MemStore results represent an upper bound on the performance of Ceph's software layers. These findings should be validated against physical storage devices, particularly high-performance \acp{SSD}, in order to allow for the construction of a complete performance model that accounts for Ceph's BlueStore---and in future Ceph releases, Crimson---overhead, write-ahead log behaviour, and storage media characteristics. Such validation would be especially useful for quantifying the performance differences between our software ceiling and achievable real-world performance, and would enable more concrete deployment recommendations for production \ac{HPC} environments.

\subsection{Expanded Fabric Evaluations}

Given the instability encountered during our benchmark, we were unable to complete Omni-Path performance testing successfully. However, the limited results we have gathered show decent scaling performance with our \ac{RPS} tuning, although we were unable to perform a full comparison across various network fabrics, and were unable to test the performance of Ceph with high-speed Ethernet networking. Therefore, full-scale testing should be expanded to cover these network fabrics, in order to evaluate whether Ceph is capable of performing well without major issues (e.g., kernel panics due to faulty networking drivers in Omni-Path). Additionally, benchmarking of the performance overhead of these fabrics relative to each other would provide an indication of the ideal network fabric for Ceph.

As the \ac{RHEL} 8 kernel vulnerability in Omni-Path networking has since been patched, revisiting the Omni-Path performance is a natural first step in order to determine whether the performance characteristics remained identical to the partial results we have gathered.

\subsection{Deployment}

The Cuttlefish in-memory cluster deployment framework developed for this evaluation demonstrated its utility in rapidly deploying rootless clusters on regular \ac{HPC} compute nodes for benchmarking purposes. However, because we only performed testing to a limited scale, future research should consider running the tool on a larger cluster in order to explore Ceph's limitations in more detail, as well as potentially applying additional limitations, such as network bandwidth or \ac{CPU} speed limitations, both of which were not possible on our cluster due to the lack of root permissions.

\subsection{\acs{RBD} Performance Evaluation}

Our evaluation focused on the \ac{RADOS} and \ac{CephFS} interfaces. However, a direct comparison with \ac{RBD}, which was the interface used in the large-scale evaluation by Clyso engineers~\cite{ceph-1tib}, would allow for the overhead of \ac{RBD} to be quantified. All interfaces ultimately operate on top of \ac{RADOS}, but \ac{CephFS} likely incurs a much larger overhead through the use of its \ac{MDS} daemons and requirement to present a fully-featured file system to clients. Such a comparison could further isolate the level of overhead introduced by Ceph's other storage services.

\subsection{Benchmarking Scripts}

As part of our testing, we have developed a tool to automatically collect scaling performance data from \ac{IOR} and mdtest into a unified SQLite database for fast exploration of multi-dimensional performance data. We believe that this tool is useful for other storage systems, and can be extended to support a wide variety of benchmarking configurations beyond only \ac{IOR} and mdtest for \ac{CephFS}, providing a more detailed view into the strengths and limitations of \ac{HPC} storage clusters under a variety of client-side scaling levels and workloads.

\subsection{Chaos Engineering}

In order to definitively validate Ceph's suitability as an \ac{HPC} file system, future evaluations should perform benchmarking on physical clusters with fault injections targeted towards Ceph's weaknesses and fault tolerance systems, such as the intentional introduction of slow \acp{OSD} to evaluate the cluster's performance, or the nondeterministic termination of nodes and/or daemons. Observation of Ceph's fault recovery performance is essential in order to prove Ceph's robustness in production environments where hardware failures occur regularly.

\subsection{\acs{CephFS} Performance Mitigations}

Given our observed performance characteristics, as well as the \ac{I/O} impact of \ac{AI} and other new workloads on the file system, performing a loopback mount of disk images~\cite{Krotkiewski_2017} could serve as a potential way to mitigate the metadata bottleneck, while retaining the benefits of the high data path performance observed in our benchmarks. Effectively, this method shifts a large amount of the \ac{I/O} workload to the highly scalable \ac{RADOS} layer, rather than the \ac{CephFS} layer's \ac{MDS} daemons. However, since we did not have the access required to perform such a test, we cannot provide the empirical data to support this hypothesis. Regardless, we consider this as a potentially interesting avenue for future research.

